\section{Introduction}

Ticket management systems record physical defects for inspection, prioritization, repair, and follow-up.
Ideally, each ticket holds one incident's history, but users often photograph the same crack more than once, creating duplicate tickets that fragment this history and cause redundant work.
Each new observation must therefore be linked to an existing ticket or open one for unrecorded damage.

This is a visual re-identification (re-ID) problem: the same crack must be recognized despite changes in viewpoint, distance, field of view, illumination, and visibility, and despite growth by extension, widening, or branching, while new damage must be distinguished from known damage.
Errors in either direction lead to duplicate tickets or misattributed history.
Crack detection and segmentation work focuses on single images.
CrackSeg9k~\cite{kulkarni2022crackseg9k} and OmniCrack30k~\cite{benz2024omnicrack30k} contain no cross-observation identities, so observations cannot be associated with tickets.
Whether person and object re-ID~\cite{zhou2019osnet,qian2024reidtrends} transfers to cracks is unclear: distinct defects may look alike while one defect can vary greatly.

This work introduces \textsc{CrackID} to evaluate whether photographs depict the same wall crack: 140 photographs of 17 walls, each photographed during a single visit.
Real photographs capture within-visit acquisition changes; photometric and growth changes are probed only synthetically.
Cross-visit evaluation and ticket-level duplicate reduction are left for future work.
All accuracy figures are point estimates from a small wall-clustered sample.

We assess ten zero-shot matchers using \textsc{CrackID}.
The protocol separates retrieval from the same-defect decision and includes open-set identification with rejection: the top-ranked match is accepted only above a validation-calibrated threshold, otherwise the query is rejected as a new ticket. Rejection is necessary because a high ranking alone is insufficient evidence for a link.
To test whether matchers recognize the crack or only wall context, we use crack-erased controls and a geometric reference that aligns wall images and compares projected crack masks.

Three research questions guide the study:
\begin{description}
  \item[\textbf{RQ1.}] Can the evaluated methods associate observations of the same defect under varying image acquisition conditions?
  \item[\textbf{RQ2.}] How do wall context and geometric alignment contribute to defect re-identification?
  \item[\textbf{RQ3.}] How do two selected approaches perform under synthetic photometric and growth changes?
\end{description}
RQ1 and RQ2 use real single-visit photographs; RQ3 is exploratory synthetic work.

The contributions are as follows:
(i) the cross-observation identities of \textsc{CrackID}, together with their annotation provenance (Section~\ref{sec:annotation});
(ii) a benchmark for retrieval, same-defect decisions, and open-set identification, addressing RQ1 (Sections~\ref{sec:protocol},~\ref{sec:baselines}, and~\ref{sec:results});
(iii) context controls, a geometric reference, and coverage analysis, addressing RQ2 (Sections~\ref{sec:dataset},~\ref{sec:method},~\ref{sec:controls}, and~\ref{sec:coverage}); and
(iv) a synthetic revisit probe for RQ3 (Section~\ref{sec:hybridbench}).
